\documentclass[sigconf,anonymous,review]{acmart}
%% DAC submissions are double-blind: keep [anonymous,review] until camera-ready.
%% DAC research papers have historically been 6 pages in this format.
%% Re-check the DAC 2027 CFP when posted and adjust DAC_PAGE_LIMIT in compile.sh.

\usepackage{placeins}
\usepackage{enumitem}
\usepackage{booktabs}
\setlist[itemize]{topsep=4pt plus 0pt minus 0pt, itemsep=3pt plus 0pt minus 0pt, parsep=0pt}

\AtBeginDocument{%
  \providecommand\BibTeX{{Bib\TeX}}}

%% Rights placeholders -- replace from the ACM rights e-mail on acceptance.
\setcopyright{acmlicensed}
\copyrightyear{2027}
\acmYear{2027}
\acmDOI{XXXXXXX.XXXXXXX}
\acmConference[DAC '27]{64th ACM/IEEE Design Automation Conference}{July 2027}{San Jose, CA, USA}
\acmISBN{978-1-4503-XXXX-X/27/07}

\begin{document}

%% TITLE -- placeholder. Settle once the headline result exists; do not
%% pre-commit the framing before the measurement.
\title{Schedule Exploration for Multicore RTL Simulation}

\begin{abstract}
%% Write this LAST, from the results. A placeholder here on purpose: an
%% abstract written before the numbers is a hypothesis, and hypotheses have
%% been wrong repeatedly on this project.
\end{abstract}

\maketitle

\section{Introduction}
%% The observation: cycle-accurate RTL simulation of a multicore is
%% deterministic, so it explores one interleaving per test program regardless of
%% how many times it is run. Silicon does not have this problem, which is why
%% post-silicon multicore validation works; simulation loses that coverage
%% source precisely pre-silicon, where bugs are cheapest to fix.
%% Evidence: docs/research/12, section 2.

\section{Background and Related Work}
%% Two axes, and they are orthogonal. Program generation: TSOtool (ISCA'04),
%% McVerSi (HPCA'16), MTraceCheck, DifuzzRTL, TheHuzz, Cascade, ProcessorFuzz.
%% Schedule exploration: CHESS (OSDI'08), PCT (ASPLOS'10), SegFuzz -- all
%% SOFTWARE. Oracles: DiffTest (multicore, commit-level -- state this
%% correctly), RTLcheck, RealityCheck.
%% Do NOT claim "first multicore fuzzer" or "unsolved oracle problem".
%% See docs/research/11 section 0 for the retracted claims.

\section{Motivation: Determinism as a Coverage Wall}
%% Fig.~\ref{fig:determinism}

\section{Schedule Exploration}
\subsection{Perturbation Mechanism}
%% One stall input per core gating a single ready signal. Tied low, the design
%% is bit-identical to the unmodified RTL.
\subsection{Soundness}
%% Delay injection cannot create a schedule the hardware could not itself
%% reach, so there are no false positives by construction. Contrast with fault
%% injection, which manufactures impossible states and requires triage.
\subsection{Scheduling Policy}
%% PCT-style bounded priority changes; bug depth at cycle granularity is a
%% genuine open question -- say so.

\section{Experimental Setup}
%% Chiron: quad-core RV64IMA OoO, AXI-ACE coherence, boots SMP Linux.
%% Verilator 4.038. Oracles: quad lockstep, swmr_probe, mt-litmus.
%% Corpus: docs/research/04.

\section{Evaluation}
\subsection{Bugs Found per Simulated Cycle}
\subsection{Comparison with Existing Practice}
%% Ablations, not unfair cross-tool claims: repeated deterministic runs;
%% ad-hoc random delay injection (UVM practice); long-run brute force (Linux
%% boot); ours. Cascade/DifuzzRTL are single-core and target other DUTs -- say
%% why a direct comparison would be misleading rather than manufacturing one.
\subsection{Interleaving Coverage}
\subsection{Sensitivity and Cost}

\section{Limitations}
%% Whole-core stalls may be too coarse for line-level races; PCT's bug-depth
%% notion needs redefining at cycle granularity; single DUT.

\section{Conclusion}

\bibliographystyle{ACM-Reference-Format}
\bibliography{refs}

\end{document}
